\section{Presentacion convertida en practica: MANOVA con respuestas fisiologicas en plantas}

Esta seccion nace de una presentacion existente y se incorpora al manual para que el material no permanezca aislado. La presentacion deja de ser un apoyo oral y se transforma en una ruta de estudio: problema, variables, datos, analisis, decision y actividad.

\begin{idea}
Decidir cuando varias respuestas correlacionadas deben analizarse como sistema.
\end{idea}

\subsection{Evidencia didactica extraida}

\begin{itemize}
\item Análisis Multivariado de la Varianza con Tres Variables Dependientes
\item Evaluar si los factores experimentales (luz, riego y fertilizante) influyen de manera
\item significativa en tres variables de respuesta fisiológica de las plantas:
\item Tipo: Factorial 3 × 3 × 3 reducido a 10 tratamientos
\end{itemize}

\subsection{Como entra al manual}

El tema se integra al capitulo relacionado con \textbf{manova}. Si el estudiante solo observa la presentacion, ve una secuencia visual. En el manual, esa misma secuencia se expande con la pregunta de ingenieria, la razon del metodo, la interpretacion y los limites de la conclusion.

\subsection{Practica asociada}

La practica generada pide al estudiante transformar la presentacion en una decision: definir variables, organizar datos, graficar, aplicar el metodo y redactar una conclusion prudente. Si no existia practica, el sistema crea una primera version editable. Si existia una practica relacionada, esta se usa para enriquecer la narrativa y no como bloque aislado.

\subsection{Actividad de evaluacion}

Explique que parte de la presentacion corresponde a contexto, cual corresponde a diseno, cual corresponde a evidencia y cual corresponde a decision. Si falta alguno de estos elementos, proponga como completarlo.
